% problem-000445 7 铜铅电解与水解平衡
\section*{7. 铜铅电解与水解平衡}
在 298.2 K， $p ^ { \ominus }$ 时，以 Cu 为阴极，⽯墨为阳极，电解含有 PbCl (0.0100 mol L−1)和 $\mathrm { C u C l _ { 2 } } \left( 0 . 0 5 0 0 \mathrm { m o l } \right.$ L−1) 的⽔溶液（不⽤盐桥），忽略电解过程中的超电势，活度因⼦均为1.00。已知 $\mathsf { J } \varphi ^ { \ominus } ( \mathrm { P b ^ { 2 + } } / \mathrm { P b } ) = - 0 . 1 2 5 \mathrm { V } , \ \varphi ^ { \ominus }$ (Cu2+/Cu) = 0.337 V， $\varphi ^ { \ominus } ( \mathrm { C l _ { 2 } } / \mathrm { C l ^ { - } } ) = 1 . 3 6 ~ \mathrm { V } ,$ $\varphi ^ { \ominus } ( \mathrm { O _ { 2 } / H _ { 2 } O } ) = 1 . 2 2 9 \mathrm { V } .$ ，⽔的 $K _ { \mathrm { w } } = 1 . 0 0 \times 1 0 ^ { - 1 4 } \mathrm { , }$ 。通过计算说明：7-1 为使第⼆种⾦属析出，⾄少需要加多⼤电压？此时第⼀种⾦属离⼦的浓度是多少？

\subsection*{7-2}
若很短时间即停⽌通电，此时 $\mathrm { C u C l _ { 2 } }$ 的浓度为 0.0400 mol L−1，反应槽 $\mathrm { p H } = 4 . 0 0 $ ，温度同上，放置⼀段时间后是否有CuOH⽣成？（不需要考虑是否进⼀步⽣成 $\mathrm { C u _ { 2 } O }$ 等）

$
\text { 已知 } K _ {\mathrm{sp}} ^ {\ominus} = 1 0 ^ {- 1 4}, \varphi^ {\ominus} (\mathrm{Cu} ^ {2 +} / \mathrm{Cu}) = 0. 1 6 0 \mathrm{V} 。
$

\subsection*{7-3}
分析纯的CuCl_{2}溶于蒸馏⽔后，加酸调 $\mathrm { p H } = 3 . 0 0$ $\mathrm { C u ^ { 2 + } }$ 仍可⽔解，此时假定只⽣成 $\mathrm { [ C u ( O H ) ] }$ +，若此⽔解反应的平衡常数 $K = 0 . 1 2 5$ ，计算1升蒸馏⽔中溶解的 $\mathrm { C u C l _ { 2 } }$ 的质量。

\subsection*{7-4}
电化学中通常将实际分解电压与理论分解电压之间的差称为超电压或超电势。 $\mathrm { O _ { 2 } ( g ) }$ 在⽯墨上析出存在超电势，若其超电势为0.800 V，通过计算说明⽤氯碱⼯业（电解饱和⻝盐⽔）法制Cl 是否可⾏？（设Cl−的浓度为0.120 mol L−1）

\subsection*{7-5}
上个世纪五⼗年代末开发的在均相溶液中⽤ $\mathrm { P d C l _ { 2 } / C u C l _ { 2 } }$ 催化剂在温和条件将⼄烯氧化制⼄醛，是过渡⾦属催化剂在⼯业上著名的应⽤实例。其总反应可表⽰为下⾯3个催化流程：

$
\begin{array}{r l}&\mathrm {PdCl_ {2}} + \mathrm {C_ {2} H_ {4}} + \mathrm {H_ {2} O} \rightarrow \mathrm {CH_ {3} CHO} + \mathrm{Pd} + 2 \mathrm{HCl}\\&\mathrm {Pd + 2CuCl_ {2}} \rightarrow \mathrm {PdCl_ {2}} + 2 \mathrm{CuCl}\\&2 \mathrm{CuCl} + 2 \mathrm{HCl} + 1 / 2 \mathrm {O_ {2}} \rightarrow 2 \mathrm {CuCl_ {2}} + \mathrm {H_ {2} O}\end{array}
$

⼄烯氧化制⼄醛的反应速率可以通过反应过程中⼄烯体积的减少或压强的降低来测定。实验发现，当反应刚开始时，⼄烯很快地被溶液吸收，且吸收量超过了饱和该溶液所需要的量，随后吸收速率减慢。在中等的 H+和氯化物浓度下，其实验速率⽅程可表⽰为：

$
r = \frac {\mathrm{d} c (\mathrm{C} _ {2} \mathrm{H} _ {4})}{\mathrm{d} t} = k c (\mathrm{PdCl} _ {4} ^ {2 -}) ^ {x} c (\mathrm{C} _ {2} \mathrm{H} _ {4}) ^ {y} c (\mathrm{H} ^ {+}) ^ {z} c (\mathrm{Cl} ^ {-}) ^ {\omega}
$

⽬前学术界公认的机理为：

$
① \left[ \mathrm{PdCl} _ {4} \right] ^ {2 -} + \mathrm{C} _ {2} \mathrm{H} _ {4} \underset {k = 1} {\overset {k _ {1}} {\rightleftharpoons}} \left[ \mathrm{PdCl} _ {3} (\mathrm{C} _ {2} \mathrm{H} _ {4}) \right] ^ {-} + \mathrm{Cl} ^ {-} \quad \text { 快反应 }
$

$
② \left[ \mathrm{PdCl} _ {3} (\mathrm{C} _ {2} \mathrm{H} _ {4}) \right] ^ {-} + \mathrm{H} _ {2} \mathrm{O} \underset {k _ {- 2}} {\overset {k _ {2}} {\rightleftharpoons}} \left[ \mathrm{PdCl} _ {2} (\mathrm{H} _ {2} \mathrm{O}) (\mathrm{C} _ {2} \mathrm{H} _ {4}) \right] + \mathrm{Cl} ^ {-} \quad \text { 快反应 }
$

$
③ \left[ \mathrm{PdCl} _ {2} (\mathrm{H} _ {2} \mathrm{O}) (\mathrm{C} _ {2} \mathrm{H} _ {4}) \right] \underset {k _ {- 3}} {\overset {k _ {3}} {\rightleftharpoons}} \left[ \mathrm{PdCl} _ {2} (\mathrm{OH}) (\mathrm{C} _ {2} \mathrm{H} _ {4}) \right] ^ {-} + \mathrm{H} ^ {+} \quad \text { 快反应 }
$

$
④ [ \mathrm{PdCl} _ {2} (\mathrm{OH}) (\mathrm{C} _ {2} \mathrm{H} _ {4}) ] ^ {-} \xrightarrow {\kappa_ {4}} [ \mathrm{Cl} _ {2} \mathrm{PdCH} _ {2} \mathrm{CH} _ {2} \mathrm{OH}) ] ^ {-}
$

$
⑤ \left[ \mathrm{Cl} _ {2} \mathrm{PdCH} _ {2} \mathrm{CH} _ {2} \mathrm{OH}) \right] ^ {-} \xrightarrow {k _ {5}} \mathrm{CH} _ {3} \mathrm{CHO} + \mathrm{Pd} + \mathrm{HCl} + \mathrm{Cl} ^ {-}
$

同位素⽰踪表明 $\mathrm { C _ { 2 } H _ { 4 } }$ 在 $\mathrm { D } _ { 2 } \mathrm { O }$ 中氧化不含 D， $\mathrm { C _ { 2 } D _ { 4 } }$ 的氧化产物只有 $\mathrm { C D _ { 3 } C D O }$ 且 $\mathrm { C _ { 2 } H _ { 4 } }$ 和 $\mathrm { C _ { 2 } D _ { 4 } }$ 的反应速率基本相同。

\subsection*{7-5}
-1 通过所给机理及相关信息，判断④和⑤哪个是决速步骤，并推导�、�、�、�的值。

\subsection*{7-5}
-2 为何在反应开始阶段，⼄烯吸收速率很快且是过量的，随后吸收减少？

\subsection*{7-5}
-3 在历程③中，经历五配位中间体，请判断产物 $\mathrm { P d C l _ { 2 } ( O H ) ( C _ { 2 } H _ { 4 } ) } ]$ −是顺式还是反式更有利于后续的反应？

\subsection*{7-5}
-4 历程⑤重排过程中，有⼀个Pd的四配位中间体，画出该中间体的⽴体结构。

\subsection*{7-5}
-5 在实际⼯业⽣产中，要求控制Cl/Pd（原⼦⽐个数）在⼀个合理的范围内，试说明两个主要的原因。

\subsection*{7-5}
-6 原料⼄烯需预先纯化除去 $\mathrm { C _ { 2 } H _ { 2 } } .$ ，试说明理由。
